\documentclass{article}
\usepackage{amsmath}
\usepackage{amssymb}
\usepackage{graphicx}
\usepackage{float}
\usepackage{subcaption}
\usepackage{array}
\usepackage[colorlinks=true, linkcolor=blue, citecolor=blue, urlcolor=blue]{hyperref}
\numberwithin{equation}{section}
\setlength{\parindent}{0pt}
\setlength{\parskip}{0pt}
\title{Beverton--Holt Model for Populations on Patches with Harvesting:\\ Preventing Extinction and Optimizing Yield}
\author{}
\date{\today}
\begin{document}
\maketitle
\begin{abstract}
This paper studies Beverton--Holt population dynamics on one patch and
on two linked patches, with harvesting represented by fishing effort.
For the single-population model, we derive the equilibria, the
extinction criterion, and the effort that maximizes sustainable yield.
For two linked patches, dispersal makes the fate of each patch depend
on the other patch, so extinction and optimal harvesting must be
studied as properties of the coupled system. We derive extinction
criteria from the linearization at the
extinction equilibrium and use the spectral radius to describe when the
linked population can persist. We then formulate the yield-maximization
problem for the linked model with fishing. A change from effort
variables to equilibrium population variables makes the problem
convex, so the Karush--Kuhn--Tucker conditions give global optimality
certificates. This leads to several possible maximizing strategies:
an interior fishing strategy, complete harvest in one patch, protection
of a source patch while harvesting a sink patch, and no-take marine
reserve boundaries. Numerical plots are used throughout to compare the
analytic conditions with simulated dynamics and to show how the optimal
strategy changes with the model parameters.
\end{abstract}
\section{The Single Population}
\subsection{Beverton--Holt Model}
We model the population with the Beverton--Holt map. This model was
introduced by Beverton and Holt in their study of exploited fish
populations \cite{beverton-holt-1957}. Let $N_t \geq 0$ denote the
population at time $t$. One intuitive way to write the model is
\begin{equation*}
    N_{t+1}
    = \frac{rN_t}{1+\frac{r-1}{K_{\mathrm{eq}}}N_t}.
\end{equation*}
Here $r$ controls the growth of a very small population, and
$K_{\mathrm{eq}}$ is the positive equilibrium when $r>1$.
Throughout this paper we use the more convenient form obtained by the
change of variables
\begin{equation*}
    K=\frac{K_{\mathrm{eq}}}{r-1}.
\end{equation*}
This gives the equivalent form
\begin{equation}
    \label{eq:bh-map}
    N_{t+1} = F(N_t)
            = \frac{K r N_t}{K + N_t},
\end{equation}
which is algebraically simpler in the later sections. In this
parametrization $K$ is the equilibrium scale scaled by a factor of
$1/(r-1)$.
The low-density meaning of $r$ is unchanged: since $F'(0)=r$, a very
small population is multiplied approximately by $r$ in one time step.\\
A cobweb plot gives a geometric picture of the iteration. Starting from
$N_t$, we apply the function by moving vertically to the point
$(N_t,F(N_t))$. The height of this point is the next population value,
$N_{t+1}$. To use it as the next input, we move horizontally to the
diagonal $N_{t+1}=N_t$, where the same number is read on the horizontal
axis. Repeating these vertical and horizontal moves shows how the
sequence approaches an equilibrium.
\begin{figure}[H]
    \centering
    \includegraphics[width=0.65\textwidth]{plots/beverton_holt_cobweb_r1_4_k1.pdf}
    \caption{Cobweb plot for the Beverton--Holt map \eqref{eq:bh-map}
    with $r=1.4$ and $K=1$. Starting from a small population, the
    iteration moves toward the positive equilibrium $N^*=K(r-1)$. The
    inset zooms in near the equilibrium.}
\end{figure}
\subsection{Equilibria}
An equilibrium $N^*$ satisfies $F(N^*)=N^*$. Using \eqref{eq:bh-map},
this means
\begin{equation*}
    N^* = \frac{K r N^*}{K + N^*}.
\end{equation*}
One solution is the extinction equilibrium
\begin{equation*}
    N^*_0=0.
\end{equation*}
For a positive equilibrium, $N^*>0$, we divide by $N^*$ and obtain
\begin{align*}
    1 &= \frac{K r}{K+N^*},\\
    K+N^* &= Kr,\\
    N^* &= K(r-1).
\end{align*}
The positive equilibrium exists only when $r>1$ since the population can
not be negative. Thus,
\begin{equation*}
    \begin{cases}
        r\leq 1: & N^*=0 \text{ is the only equilibrium},\\
        r>1: & N^*=0 \text{ and } N^*=K(r-1) \text{ are equilibria}.
    \end{cases}
\end{equation*}
\subsection{Stability}
Stability describes what happens after a small perturbation. If the
population starts close to an equilibrium and stays close, the
equilibrium is stable. If the population also returns to the equilibrium
as time goes on, the equilibrium is asymptotically stable.

Formally, for a discrete map $N_{t+1}=F(N_t)$, an equilibrium $N^*$ is
stable if for every $\varepsilon>0$ there exists $\delta>0$ such that
\begin{equation*}
    |N_0-N^*|<\delta
    \quad \Longrightarrow \quad
    |N_t-N^*|<\varepsilon
    \quad \text{for all } t\geq 0.
\end{equation*}
It is locally asymptotically stable if, in addition, all sufficiently
nearby initial populations converge to $N^*$:
\begin{equation*}
    |N_0-N^*|<\delta
    \quad \Longrightarrow \quad
    \lim_{t\to\infty}N_t=N^*.
\end{equation*}
For the one-dimensional maps used in this paper, stability can be
checked from the derivative. If $F$ is continuously differentiable near
an equilibrium $N^*$ and
\begin{equation*}
    |F'(N^*)|<1,
\end{equation*}
then $N^*$ is locally asymptotically stable. If $|F'(N^*)|>1$, then
$N^*$ is unstable. The case $|F'(N^*)|=1$ must be checked separately.
This is a special case of the Banach fixed-point theorem
\cite[pp.~7--22]{kinderlehrer-stampacchia-1980}. To avoid using this
heavier theorem, we prove the one-dimensional stable case directly.
Choose a number $q$ such that
\begin{equation*}
    |F'(N^*)|<q<1.
\end{equation*}
Since $F'$ is continuous, there is a small interval around $N^*$ where
$|F'(N)|\leq q$. For $N$ in this interval, the mean value theorem gives
\begin{equation*}
    |F(N)-F(N^*)|\leq q|N-N^*|.
\end{equation*}
Because $F(N^*)=N^*$, this becomes
\begin{equation}
    \label{eq:local-contraction}
    |F(N)-N^*|\leq q|N-N^*|.
\end{equation}
Now apply this estimate to the iterates. First,
\begin{equation*}
    |N_1-N^*|
    = |F(N_0)-N^*|
    \leq q|N_0-N^*|.
\end{equation*}
Then, using the same estimate again,
\begin{equation*}
    |N_2-N^*|
    = |F(N_1)-N^*|
    \leq q|N_1-N^*|
    \leq q^2|N_0-N^*|.
\end{equation*}
Repeating this argument gives
\begin{equation*}
    |N_t-N^*|\leq q^t|N_0-N^*|.
\end{equation*}
Since $0<q<1$, we have $q^t\to 0$, and therefore $N_t\to N^*$.

The derivative of the population map \eqref{eq:bh-map} is
\begin{equation*}
    F'(N)=\frac{K^2r}{(K+N)^2}.
\end{equation*}
At the extinction equilibrium,
\begin{equation*}
    F'(0)=r.
\end{equation*}
This derivative measures the growth of a very small population. Thus,
for $r<1$, extinction is locally asymptotically stable by the derivative
test. For $r>1$, extinction is unstable. The threshold case $r=1$ is not
decided by the derivative test, but it can be checked directly:
\begin{equation*}
    F(N)
    =\frac{KN}{K+N}
    =N\frac{K}{K+N}
    <N
    \quad \text{for every } N>0.
\end{equation*}
Thus positive populations decrease. Since the sequence is nonnegative,
it has a limit $L\geq 0$. By continuity, this limit must satisfy
$F(L)=L$. At $r=1$, the only nonnegative equilibrium is $L=0$, so
positive populations decrease toward extinction. The extinction
criterion for the model without fishing is therefore
\begin{equation*}
    r \leq 1.
\end{equation*}
At the positive equilibrium,
\begin{equation*}
    F'\bigl(K(r-1)\bigr)=\frac{1}{r}.
\end{equation*}
Consequently, whenever the positive equilibrium exists, $r>1$, it is
locally asymptotically stable.
\section{Introduction of Fishing Effort}
We now introduce a fishing effort parameter $e$, where
\begin{equation*}
    0 \leq e \leq 1.
\end{equation*}
The value $e=0$ means no fishing, while $e=1$ means complete removal
after reproduction. With fishing, the one-population model becomes
\begin{equation}
    \label{eq:fishing-map}
    N_{t+1}
    = F_e(N_t)
    = (1-e)\frac{K r N_t}{K+N_t}.
\end{equation}
\subsection{Equilibria}
An equilibrium satisfies $F_e(N^*)=N^*$. Using \eqref{eq:fishing-map},
this gives
\begin{equation*}
    N^* = (1-e)\frac{K r N^*}{K+N^*}.
\end{equation*}
Again, $N^*=0$ is always an equilibrium. For a positive equilibrium,
$N^*>0$, we divide by $N^*$ and obtain
\begin{align*}
    1 &= (1-e)\frac{K r}{K+N^*},\\
    K+N^* &= K r(1-e),\\
    N^* &= K\bigl(r(1-e)-1\bigr).
\end{align*}
Therefore, the positive equilibrium exists only when
\begin{equation*}
    r(1-e)>1.
\end{equation*}
Equivalently, if $r>1$, persistence requires
\begin{equation*}
    e < 1-\frac{1}{r}.
\end{equation*}
If this condition fails, then extinction is the only nonnegative
equilibrium.
The derivative of the fishing map \eqref{eq:fishing-map} is
\begin{equation*}
    F_e'(N)=\frac{K^2r(1-e)}{(K+N)^2}.
\end{equation*}
At extinction,
\begin{equation*}
    F_e'(0)=r(1-e).
\end{equation*}
Therefore, the extinction criterion with fishing is
\begin{equation*}
    r(1-e)\leq 1.
\end{equation*}
Equivalently, when $r>1$, extinction occurs when
\begin{equation}
    \label{eq:fishing-extinction-effort}
    e \geq 1-\frac{1}{r}.
\end{equation}
When $r(1-e)<1$, extinction is locally stable. When $r(1-e)>1$,
extinction is unstable and the positive equilibrium exists.
At the positive equilibrium,
\begin{equation*}
    F_e'\bigl(K(r(1-e)-1)\bigr)=\frac{1}{r(1-e)}.
\end{equation*}
Whenever the positive equilibrium exists, $r(1-e)>1$, this derivative
is less than $1$, so the positive equilibrium is locally stable.
\begin{figure}[H]
    \centering
    \includegraphics[width=0.8\textwidth]{plots/fishing_equilibrium_r1_4_k1.pdf}
    \caption{Equilibrium population as a function of fishing effort for
    $r=1.4$ and $K=1$. The dashed vertical line marks the extinction
    criterion $e_c=1-\frac{1}{r}$ from
    \eqref{eq:fishing-extinction-effort}.}
\end{figure}
\subsection{Maximizing Fishing Yield}
From a management point of view, survival alone is not the only
question. Once fishing is allowed, there is a natural pressure to ask how
much can be taken now while still leaving enough population for the same
process to continue in the future. This is the
optimization problem behind maximum sustainable yield.

Fishing effort has two competing effects. Increasing $e$ means removing
a larger fraction of the population after reproduction, but it also
lowers the equilibrium population that remains for future time steps.
Thus very small effort gives little harvest, while very large effort can
push the system to extinction. This suggests that an intermediate effort
may maximize the long-term yield.

We define the equilibrium fishing yield as the amount removed by fishing
at equilibrium. Since fishing removes the fraction $e$ after
reproduction, the yield is
\begin{equation*}
    Y(e)=eF(N^*).
\end{equation*}
At equilibrium, $N^*=(1-e)F(N^*)$, and therefore
\begin{equation*}
    F(N^*)=\frac{N^*}{1-e}.
\end{equation*}
Thus, for a positive equilibrium,
\begin{equation}
    \label{eq:yield-expression}
    \begin{aligned}
        Y(e)
            &= \frac{eN^*}{1-e}\\
            &= \frac{eK(r(1-e)-1)}{1-e}\\
            &= Ke\left(r-\frac{1}{1-e}\right).
    \end{aligned}
\end{equation}
The yield expression \eqref{eq:yield-expression} is meaningful only
below the extinction threshold in \eqref{eq:fishing-extinction-effort}.
To maximize yield, differentiate:
\begin{equation*}
    Y'(e)=K\left(r-\frac{1}{(1-e)^2}\right).
\end{equation*}
The critical point satisfies
\begin{equation*}
    r=\frac{1}{(1-e)^2}.
\end{equation*}
Therefore, for $r>1$, the yield-maximizing effort is
\begin{equation}
    \label{eq:yield-optimal-effort}
    e^*=1-\frac{1}{\sqrt{r}}.
\end{equation}
For $r=1.4$ and $K=1$, \eqref{eq:yield-optimal-effort} gives
$e^*\approx 0.155$, while the extinction threshold from
\eqref{eq:fishing-extinction-effort} is $e_c\approx 0.286$.
\begin{figure}[H]
    \centering
    \includegraphics[width=0.8\textwidth]{plots/fishing_yield_r1_4_k1.pdf}
    \caption{Equilibrium fishing yield as a function of effort for
    $r=1.4$ and $K=1$. The green dashed line marks the yield-maximizing
    effort from \eqref{eq:yield-optimal-effort}, while the red dashed
    line marks the extinction threshold from
    \eqref{eq:fishing-extinction-effort}.}
\end{figure}
\section{Linked Two-Population Model}
We now introduce two linked patches, denoted by $A$ and $B$. The
population sizes in these patches at time $t$ are $N_A(t)$ and
$N_B(t)$. Since the patches may have different local growth conditions,
we use patch-specific parameters: patch $A$ has parameters $r_A$ and
$K_A$, while patch $B$ has parameters $r_B$ and $K_B$.
The local Beverton-Holt growth functions are
\begin{equation*}
    F_A(N_A)=\frac{K_A r_A N_A}{K_A+N_A},
    \qquad
    F_B(N_B)=\frac{K_B r_B N_B}{K_B+N_B}.
\end{equation*}
The parameter $m$ controls symmetric movement, or dispersal, between the
two patches. A fraction $m$ of the post-growth population moves to the
other patch, while a fraction $1-m$ remains in the same patch.
Thus, for $0\leq m\leq 1$, the linked model without fishing is
\begin{equation}
    \label{eq:linked-model}
    \begin{aligned}
        N_A(t+1)
            &= (1-m)\frac{K_A r_A N_A(t)}{K_A+N_A(t)}
               + m\frac{K_B r_B N_B(t)}{K_B+N_B(t)},\\
        N_B(t+1)
            &= m\frac{K_A r_A N_A(t)}{K_A+N_A(t)}
               + (1-m)\frac{K_B r_B N_B(t)}{K_B+N_B(t)}.
    \end{aligned}
\end{equation}
When $m=0$, the two populations are independent. When $m=1$, the two
post-growth populations exchange locations completely at each time step.
\subsection{Equilibria}
An equilibrium of the linked model is a pair $(N_A^*,N_B^*)$ satisfying
\begin{equation*}
    \begin{aligned}
        N_A^*
            &= (1-m)\frac{K_A r_A N_A^*}{K_A+N_A^*}
               + m\frac{K_B r_B N_B^*}{K_B+N_B^*},\\
        N_B^*
            &= m\frac{K_A r_A N_A^*}{K_A+N_A^*}
               + (1-m)\frac{K_B r_B N_B^*}{K_B+N_B^*}.
    \end{aligned}
\end{equation*}
The extinction equilibrium
\begin{equation*}
    (N_A^*,N_B^*)=(0,0)
\end{equation*}
always exists.
For $m=0$, the two populations are independent, so the positive
equilibrium is simply
\begin{equation*}
    N_A^*=K_A(r_A-1),
    \qquad
    N_B^*=K_B(r_B-1),
\end{equation*}
when $r_A>1$ and $r_B>1$.
\subsection{Connect or Not to Connect}
A natural question is whether dispersal between two patches helps or hurts
the total equilibrium population. One might expect dispersal to only
redistribute the same total population between the two locations.
However, dispersal can change the total equilibrium size.

This subsection reproduces an effect from \textbf{[paper
placeholder]}. %TODO REPLACE palceholder

Consider the symmetric scale case
\begin{equation*}
    K_A=K_B=1,
    \qquad
    r_A=1.8,
    \qquad
    r_B=1.2.
\end{equation*}
Without dispersal, $m=0$, the two populations are independent, so
\begin{equation*}
    N_A^*=r_A-1=0.8,
    \qquad
    N_B^*=r_B-1=0.2.
\end{equation*}
Thus the total equilibrium population is
\begin{equation*}
    N_A^*+N_B^*=1.0.
\end{equation*}
If we allow dispersal with $m=0.2$, numerical solution of the
coupled equilibrium equations gives approximately
\begin{equation*}
    N_A^*\approx0.6259,
    \qquad
    N_B^*\approx0.4248.
\end{equation*}
Therefore,
\begin{equation*}
    N_A^*+N_B^*\approx1.0507.
\end{equation*}
In this example, dispersal increases the total equilibrium population.
\begin{figure}[H]
    \centering
    \includegraphics[width=0.8\textwidth]{plots/linked_connection_effect_ra1_8_rb1_2.pdf}
    \caption{Equilibrium populations as the movement rate $m$ varies for
    $K_A=K_B=1$, $r_A=1.8$, and $r_B=1.2$. The dashed horizontal line is
    the total equilibrium population when $m=0$. For small positive
    dispersal rates, the total equilibrium rises above the no-dispersal
    value.}
\end{figure}
\subsection{Exact Algebraic Form}
For general $m$, the two equilibrium equations are coupled. After
clearing denominators and eliminating $N_B^*$, the nonzero equilibria are
obtained exactly from a cubic polynomial
\begin{equation}
    \label{eq:linked-equilibrium-cubic}
    P(N_A^*)=0.
\end{equation}
Once a root $N_A^*$ of \eqref{eq:linked-equilibrium-cubic} is chosen,
the corresponding value of $N_B^*$ is
\begin{equation}
    \label{eq:linked-equilibrium-nb-from-na}
    N_B^*
    =
    \frac{
        K_B N_A^*
        \left(K_A m r_A - K_A r_A + K_A + N_A^*\right)
    }{
        K_AK_Bmr_B
        -K_A N_A^*mr_A
        +K_A N_A^*r_A
        -K_A N_A^*
        +K_B N_A^*mr_B
        -(N_A^*)^2
    }.
\end{equation}
Writing
$$
    P(N_A)=a_3N_A^3+a_2N_A^2+a_1N_A+a_0,
$$
the coefficients are
\begin{align*}
    a_3 &= m-1,\\
    a_2
        &= K_A m^2r_A -4K_A mr_A +2K_A m +2K_A r_A -2K_A\\
        &\qquad{} -K_Bm^2r_B +K_Bmr_B-K_Bm,\\
    a_1
        &= -2K_A^2m^2r_A^2+K_A^2m^2r_A+3K_A^2mr_A^2
           -4K_A^2mr_A+K_A^2m\\
        &\qquad{} -K_A^2r_A^2+2K_A^2r_A-K_A^2
           +2K_AK_Bm^2r_Ar_B\\
        &\qquad{} -K_AK_Bm^2r_A-2K_AK_Bm^2r_B
           -K_AK_Bmr_Ar_B\\
        &\qquad{} +K_AK_Bmr_A+2K_AK_Bmr_B-2K_AK_Bm,\\
    a_0
        &= 2K_A^2K_Bm^2r_Ar_B-K_A^2K_Bm^2r_A
           -K_A^2K_Bm^2r_B\\
        &\qquad{} -K_A^2K_Bmr_Ar_B+K_A^2K_Bmr_A
           +K_A^2K_Bmr_B-K_A^2K_Bm.
\end{align*}
Because \eqref{eq:linked-equilibrium-cubic} is cubic, its roots have a
closed form. In practice, the cubic formula is too long to be useful, so
we treat \eqref{eq:linked-equilibrium-cubic} and
\eqref{eq:linked-equilibrium-nb-from-na} as the exact algebraic
description. The biologically meaningful equilibrium is the root giving
$N_A^*>0$ and $N_B^*>0$, if such a root exists.
\subsection{Extinction Without Fishing}
To study extinction, we linearize the linked model \eqref{eq:linked-model}
at the extinction equilibrium $(0,0)$. For population $A$, the local
growth function is
\begin{equation*}
    F_A(N_A)=\frac{K_A r_A N_A}{K_A+N_A}.
\end{equation*}
Its derivative is
\begin{equation*}
    F_A'(N_A)=\frac{K_A^2 r_A}{(K_A+N_A)^2},
\end{equation*}
and therefore
\begin{equation*}
    F_A'(0)=r_A.
\end{equation*}
Similarly,
\begin{equation*}
    F_B'(0)=r_B.
\end{equation*}
Near extinction, both populations are small, so the first-order Taylor
approximations are
\begin{equation*}
    F_A(N_A)\approx r_A N_A,
    \qquad
    F_B(N_B)\approx r_B N_B.
\end{equation*}
Substituting these approximations into the linked model
\eqref{eq:linked-model} gives
\begin{equation*}
    \begin{aligned}
        N_A(t+1)
            &\approx (1-m)r_A N_A(t)+mr_BN_B(t),\\
        N_B(t+1)
            &\approx mr_A N_A(t)+(1-m)r_BN_B(t).
    \end{aligned}
\end{equation*}
Thus, the scale parameters $K_A$ and $K_B$ do not appear in the
linearized extinction dynamics. In matrix form, the linearized system is
\begin{equation*}
    \begin{pmatrix}
        N_A(t+1)\\
        N_B(t+1)
    \end{pmatrix}
    =
    M
    \begin{pmatrix}
        N_A(t)\\
        N_B(t)
    \end{pmatrix},
\end{equation*}
where
\begin{equation}
    \label{eq:linked-linear-matrix}
    M=
    \begin{pmatrix}
        r_A(1-m) & r_Bm\\
        r_Am & r_B(1-m)
    \end{pmatrix}.
\end{equation}
For a multidimensional discrete map, the derivative test is replaced by
the eigenvalue test for the Jacobian matrix: the equilibrium is locally
asymptotically stable when all eigenvalues of the linearization have
absolute value less than $1$. A full proof of this stability criterion
can be found in Elaydi \cite{elaydi-2005-stability}. Therefore, the
extinction equilibrium is locally stable when the spectral radius of the
matrix in \eqref{eq:linked-linear-matrix} satisfies
\begin{equation}
    \label{eq:linked-spectral-criterion}
    \rho(M)<1.
\end{equation}
Because $M$ is nonnegative, its spectral radius is the larger eigenvalue.
For the matrix in \eqref{eq:linked-linear-matrix}, this gives the closed
form
\begin{equation}
    \label{eq:linked-spectral-radius-closed-form}
    \rho(M)
    =
    \frac{
        (1-m)(r_A+r_B)
        +
        \sqrt{(1-m)^2(r_A-r_B)^2+4m^2r_Ar_B}
    }{2}.
\end{equation}
Thus extinction without fishing occurs exactly when
\begin{equation}
    \label{eq:linked-extinction-closed-form}
    (1-m)(r_A+r_B)
    +
    \sqrt{(1-m)^2(r_A-r_B)^2+4m^2r_Ar_B}
    <2.
\end{equation}
This recovers the independent-population condition $r_A<1$ and $r_B<1$
when $m=0$. When $m=1$, the condition becomes $r_Ar_B<1$.
\begin{figure}[H]
    \centering
    \includegraphics[width=0.75\textwidth]{plots/linked_no_fishing_extinction_m0_2.pdf}
    \caption{Extinction region for the linked two-population model
    without fishing when $m=0.2$. The blue region satisfies
    $\rho(M)<1$, and the black curve is the boundary $\rho(M)=1$.
    Random simulated points are green when the population survives and
    red when it goes extinct.}
\end{figure}
\section{Linked Two-Population Model with Fishing Effort}
We now add fishing to the linked two-population model. The effort
$e_A$ is applied to population $A$, and $e_B$ is applied to population
$B$, after local growth and movement. Thus, for $0\leq e_A,e_B\leq1$,
the model becomes
\begin{equation}
    \label{eq:linked-fishing-model}
    \begin{aligned}
        N_A(t+1)
            &= (1-e_A)
            \left(
                (1-m)\frac{K_A r_A N_A(t)}{K_A+N_A(t)}
                + m\frac{K_B r_B N_B(t)}{K_B+N_B(t)}
            \right),\\
        N_B(t+1)
            &= (1-e_B)
            \left(
                m\frac{K_A r_A N_A(t)}{K_A+N_A(t)}
                + (1-m)\frac{K_B r_B N_B(t)}{K_B+N_B(t)}
            \right).
    \end{aligned}
\end{equation}
When $e_A=e_B=0$, this reduces to the linked model
\eqref{eq:linked-model}.
\subsection{Equilibria}
An equilibrium of the linked fishing model is a pair
$(N_A^*,N_B^*)$ satisfying
\begin{equation}
    \label{eq:linked-fishing-equilibrium}
    \begin{aligned}
        N_A^*
            &= (1-e_A)
            \left(
                (1-m)\frac{K_A r_A N_A^*}{K_A+N_A^*}
                +m\frac{K_B r_B N_B^*}{K_B+N_B^*}
            \right),\\
        N_B^*
            &= (1-e_B)
            \left(
                m\frac{K_A r_A N_A^*}{K_A+N_A^*}
                +(1-m)\frac{K_B r_B N_B^*}{K_B+N_B^*}
            \right).
    \end{aligned}
\end{equation}
The extinction equilibrium
\begin{equation*}
    (N_A^*,N_B^*)=(0,0)
\end{equation*}
always exists.

When $m=0$, the two populations are independent and fishing simply
reduces the effective growth rates. The positive equilibrium is
\begin{equation*}
    N_A^*=K_A(r_A(1-e_A)-1),
    \qquad
    N_B^*=K_B(r_B(1-e_B)-1),
\end{equation*}
provided
\begin{equation*}
    r_A(1-e_A)>1,
    \qquad
    r_B(1-e_B)>1.
\end{equation*}

For general $m$, the equilibrium equations are still coupled. Solving
the first equation in \eqref{eq:linked-fishing-equilibrium} for
$N_B^*$ gives
\begin{equation}
    \label{eq:linked-fishing-nb-from-na}
    N_B^*
    =
    \frac{
        K_BN_A^*
        \left(K_A(r_A(1-e_A)(1-m)-1)-N_A^*\right)
    }{
        (N_A^*)^2
        +K_AN_A^*\left(1-r_A(1-e_A)(1-m)\right)
        -K_Bmr_B(1-e_A)(K_A+N_A^*)
    }.
\end{equation}
Substituting \eqref{eq:linked-fishing-nb-from-na} into the second
equilibrium equation gives a cubic equation for $N_A^*$, similar to the
no-fishing cubic in \eqref{eq:linked-equilibrium-cubic}. Thus fishing
changes the coefficients of the algebraic equation, but it does not
increase the degree of the equilibrium problem.
\subsection{Extinction Criterion}
To study extinction with fishing, we linearize
\eqref{eq:linked-fishing-model} at $(0,0)$. The linearized matrix is
\begin{equation}
    \label{eq:linked-fishing-linear-matrix}
    M_f=
    \begin{pmatrix}
        r_A(1-m)(1-e_A) & r_Bm(1-e_A)\\
        r_Am(1-e_B) & r_B(1-m)(1-e_B)
    \end{pmatrix}.
\end{equation}
As before, the parameters $K_A$ and $K_B$ cancel from the linearized
extinction dynamics. Extinction is determined by
\begin{equation}
    \label{eq:linked-fishing-spectral-criterion}
    \rho(M_f)<1.
\end{equation}
The boundary between extinction and persistence is given by
\begin{equation*}
    \det(I-M_f)=0.
\end{equation*}
For this matrix,
\begin{equation}
    \label{eq:linked-fishing-det-boundary}
    \det(I-M_f)=(Ce_A+D)e_B-(Ae_A+B),
\end{equation}
where
\begin{equation*}
    \begin{aligned}
        A &= -r_A + mr_A + r_Ar_B - 2mr_Ar_B,\\
        B &= -1 + r_A + r_B - mr_A - mr_B - r_Ar_B + 2mr_Ar_B,\\
        C &= r_Ar_B - 2mr_Ar_B,\\
        D &= r_B - mr_B - r_Ar_B + 2mr_Br_A.
    \end{aligned}
\end{equation*}
The threshold curve is obtained from the boundary
\begin{equation}
    \label{eq:linked-fishing-det-condition}
    (Ce_A+D)e_B-(Ae_A+B)\geq0.
\end{equation}
For a fixed effort $e_A$, the threshold form
$e_B\geq(Ae_A+B)/(Ce_A+D)$ is meaningful only on the branch where
$Ce_A+D>0$. To make the criterion valid on the whole unit square
$0\leq e_A,e_B\leq1$, we write both parts:
\begin{equation}
    \label{eq:linked-fishing-eb-threshold}
    e_B\geq\frac{Ae_A+B}{Ce_A+D}
    \qquad\text{and}\qquad
    Ce_A+D>0.
\end{equation}
The second condition does not describe another extinction mechanism. It
only restricts the threshold inequality to one side of the vertical line
$Ce_A+D=0$ where division by $Ce_A+D$ gives the correct extinction
branch. This gives the expected direction: larger fishing effort $e_B$
pushes the system toward extinction. The final extinction region is the
set of efforts for which the spectral condition
\eqref{eq:linked-fishing-spectral-criterion} holds.
The sign of $Ce_A+D$ can be understood directly from
\begin{equation*}
    Ce_A+D
    =
    r_B\left(1-m-r_A(1-2m)(1-e_A)\right).
\end{equation*}
If $1/2\leq m<1$, then $1-m>0$ and $1-2m\leq0$. Since
$r_A>0$ and $1-e_A\geq0$ on the unit square, the term
$-r_A(1-2m)(1-e_A)$ is nonnegative. Therefore
\begin{equation*}
    1-m-r_A(1-2m)(1-e_A)>0,
\end{equation*}
and hence $Ce_A+D>0$ for every $0\leq e_A\leq1$. In this case the
branch condition is automatic on the whole unit square.

If $0\leq m<1/2$, the vertical line $Ce_A+D=0$ can intersect the unit square.
In this case we need the second condition in
\eqref{eq:linked-fishing-eb-threshold}. To see why, suppose instead that
$Ce_A+D\leq0$. Then
\begin{equation*}
    r_A(1-e_A)\geq\frac{1-m}{1-2m}.
\end{equation*}
From \eqref{eq:linked-fishing-linear-matrix}, the first diagonal entry of
$M_f$ is
\begin{equation*}
    r_A(1-m)(1-e_A),
\end{equation*}
and therefore
\begin{equation*}
    r_A(1-m)(1-e_A)
    \geq
    \frac{(1-m)^2}{1-2m}
    =
    1+\frac{m^2}{1-2m}
    \geq1.
\end{equation*}
Since $M_f$ is a nonnegative matrix, its spectral radius is at least as
large as any diagonal entry. Hence $\rho(M_f)\geq1$, so extinction
cannot occur when $Ce_A+D\leq0$. Thus, for $0\leq m<1/2$, the condition
$Ce_A+D>0$ is not just a technical detail of the division; it removes a
part of the unit square where extinction is impossible.
\begin{figure}[H]
    \centering
    \begin{subfigure}{0.36\textwidth}
        \centering
        \includegraphics[width=\textwidth]{plots/linked_extinction_map_ra1_4_rb0_8_m0_3.pdf}
        \caption{$r_A=1.4$, $r_B=0.8$, $m=0.3$}
    \end{subfigure}
    \hfill
    \begin{subfigure}{0.36\textwidth}
        \centering
        \includegraphics[width=\textwidth]{plots/linked_extinction_map_ra2_0_rb1_2_m0_1.pdf}
        \caption{$r_A=2.0$, $r_B=1.2$, $m=0.1$}
    \end{subfigure}

    \begin{subfigure}{0.36\textwidth}
        \centering
        \includegraphics[width=\textwidth]{plots/linked_extinction_map_ra1_4_rb1_2_m0_1.pdf}
        \caption{$r_A=1.4$, $r_B=1.2$, $m=0.1$}
    \end{subfigure}
    \hfill
    \begin{subfigure}{0.36\textwidth}
        \centering
        \includegraphics[width=\textwidth]{plots/linked_extinction_map_ra1_4_rb1_2_m0_2.pdf}
        \caption{$r_A=1.4$, $r_B=1.2$, $m=0.2$}
    \end{subfigure}

    \begin{subfigure}{0.36\textwidth}
        \centering
        \includegraphics[width=\textwidth]{plots/linked_extinction_map_ra1_4_rb1_2_m0_3.pdf}
        \caption{$r_A=1.4$, $r_B=1.2$, $m=0.3$}
    \end{subfigure}
    \hfill
    \begin{subfigure}{0.36\textwidth}
        \centering
        \includegraphics[width=\textwidth]{plots/linked_extinction_map_ra1_4_rb1_2_m1_0.pdf}
        \caption{$r_A=1.4$, $r_B=1.2$, $m=1$}
    \end{subfigure}
    \caption{Extinction maps for the linked two-population model with
    fishing. The blue background is the theoretical extinction region
    predicted by $\rho(M_f)<1$, and the black curve is the boundary
    $\rho(M_f)=1$. Random simulated points are green when the population
    survives and red when it goes extinct.}
\end{figure}
These plots show how dispersal affects whether complete fishing in one
patch is possible. For the sequence with $r_A=1.4$ and $r_B=1.2$,
panel (c), where $m=0.1$, allows either patch to be fished completely,
provided the other patch is only lightly fished. In panel (d), where
$m=0.2$, complete fishing in patch $B$ is possible when patch $A$ is
lightly fished, but complete fishing in patch $A$ is not possible even
when patch $B$ is unfished. In panel (e), where $m=0.3$, neither patch
can be fished completely.
\subsection{Maximizing Fishing Yield}

As in the single-patch model, survival is not the only question. Once the
linked system can persist, we can ask which fishing efforts give the
largest sustainable equilibrium yield. The difference is that now the two
patches interact through dispersal, so changing the effort in one patch
also changes the equilibrium population in the other. This section
formulates that optimization problem and then classifies the possible
maximizers.

\subsubsection{Karush--Kuhn--Tucker Conditions}

We will actively use the Karush--Kuhn--Tucker conditions to certify the
maximizing candidate in each case. Since the yield problem is a
maximization problem, we state the conditions explicitly in the
maximization convention. Consider
\begin{equation}
    \label{eq:kkt-general-problem}
    \text{maximize } f(x)
    \qquad\text{subject to}\qquad
    g_i(x)\leq0,\quad i=1,\ldots,s.
\end{equation}
A constraint is called \textbf{active} at $x^*$ if $g_i(x^*)=0$;
otherwise it is \textbf{inactive}. In the active-set language, each case
below corresponds to a choice of which constraints are active.

For this maximization convention, define the Lagrangian by
\begin{equation}
    \label{eq:kkt-lagrangian}
    \mathcal L(x,\lambda)
    =
    f(x)-\sum_{i=1}^s \lambda_i g_i(x).
\end{equation}
The minus sign matches the convention $g_i(x)\leq0$ for a maximization
problem.

The KKT conditions are the following. First, the point must satisfy the
original constraints:
\begin{equation}
    \label{eq:kkt-primal-feasibility}
    g_i(x^*)\leq0,
    \qquad i=1,\ldots,s.
\end{equation}
This is called \textbf{primal feasibility}.

Second, every inequality multiplier must be nonnegative:
\begin{equation}
    \label{eq:kkt-dual-feasibility}
    \lambda_i\geq0,
    \qquad i=1,\ldots,s.
\end{equation}
This is called \textbf{dual feasibility}.

Third, each multiplier can be nonzero only when its constraint is active:
\begin{equation}
    \label{eq:kkt-complementary-slackness}
    \lambda_i g_i(x^*)=0,
    \qquad i=1,\ldots,s.
\end{equation}
This is called \textbf{complementary slackness}.

Finally, the derivative of the Lagrangian with respect to $x$ must vanish:
\begin{equation}
    \label{eq:kkt-stationarity}
    \nabla_x\mathcal L(x^*,\lambda)
    =
    \nabla f(x^*)-\sum_{i=1}^s \lambda_i \nabla g_i(x^*)=0.
\end{equation}
This is the \textbf{stationarity condition}.

For a concave maximization problem over a convex feasible set, these
conditions are sufficient for global optimality; see Boyd and Vandenberghe
\cite[Chapter~5]{boyd-vandenberghe-2004}. Thus, once we show that a
candidate satisfies
\eqref{eq:kkt-primal-feasibility}--\eqref{eq:kkt-stationarity},
we have proved that it is a global maximizer, not merely a local candidate.

\subsubsection{Problem Formulation}

We first formulate the problem in the original effort variables. For fixed
parameters $K_A,K_B>0$, $r_A,r_B\geq0$, and $0\leq m\leq1$, define
\begin{equation}
    \label{eq:growth-functions}
    F_A(N_A)=\frac{K_A r_A N_A}{K_A+N_A},
    \qquad
    F_B(N_B)=\frac{K_B r_B N_B}{K_B+N_B}.
\end{equation}
After growth and dispersal, but before harvesting, the patch populations are
\begin{equation}
    \label{eq:preharvest-growth}
    \begin{aligned}
    G_A(N_A,N_B)&=(1-m)F_A(N_A)+mF_B(N_B),\\
    G_B(N_A,N_B)&=mF_A(N_A)+(1-m)F_B(N_B).
    \end{aligned}
\end{equation}
For efforts $0\leq e_A,e_B\leq1$, an equilibrium satisfies
\begin{equation}
    \label{eq:effort-equilibrium}
    N_A=(1-e_A)G_A(N_A,N_B),
    \qquad
    N_B=(1-e_B)G_B(N_A,N_B).
\end{equation}
If the positive equilibrium determined by $(e_A,e_B)$ is denoted by
$N_A^*(e_A,e_B)$ and $N_B^*(e_A,e_B)$, then the equilibrium yield is
\begin{equation}
    \label{eq:effort-yield}
    Y(e_A,e_B)
    =
    e_A G_A(N_A^*,N_B^*)
    +
    e_B G_B(N_A^*,N_B^*).
\end{equation}
The effort formulation of the maximization problem is therefore
\begin{equation}
    \label{eq:effort-maximization-problem}
    \max_{0\leq e_A,e_B\leq1} Y(e_A,e_B),
\end{equation}
where the populations $N_A^*(e_A,e_B)$ and $N_B^*(e_A,e_B)$ are determined
implicitly by \eqref{eq:effort-equilibrium}.

This is the biologically direct formulation, but it is not the most
convenient one for calculation. The difficulty is that the variables being
optimized, $e_A$ and $e_B$, also change the equilibrium populations. To make
the constraints explicit, we reformulate the same problem using the
equilibrium populations $(N_A,N_B)$ as the variables.

From \eqref{eq:effort-equilibrium}, whenever $G_A,G_B>0$ the efforts
that maintain a chosen equilibrium pair $(N_A,N_B)$ are
\begin{equation}
    \label{eq:effort-recovery}
    e_A=1-\frac{N_A}{G_A(N_A,N_B)},
    \qquad
    e_B=1-\frac{N_B}{G_B(N_A,N_B)}.
\end{equation}
The restrictions $0\leq e_A,e_B\leq1$ are therefore equivalent to
\begin{equation}
    \label{eq:population-domain-inequalities}
    0\leq N_A\leq G_A(N_A,N_B),
    \qquad
    0\leq N_B\leq G_B(N_A,N_B).
\end{equation}
Thus the feasible population domain is
\begin{equation}
    \label{eq:population-domain}
    \mathcal D
    =
    \left\{
    (N_A,N_B)\in\mathbb R_{\geq0}^2:
    N_A\leq G_A(N_A,N_B),\;
    N_B\leq G_B(N_A,N_B)
    \right\}.
\end{equation}
In this domain, the yield becomes
\begin{equation}
    \label{eq:population-yield-before-separation}
    Y(N_A,N_B)
    =
    G_A(N_A,N_B)-N_A
    +
    G_B(N_A,N_B)-N_B.
\end{equation}
Since dispersal only redistributes the growth terms,
$G_A+G_B=F_A+F_B$, and hence
\begin{equation}
    \label{eq:population-yield}
    Y(N_A,N_B)
    =
    F_A(N_A)-N_A
    +
    F_B(N_B)-N_B.
\end{equation}
The population-domain maximization problem is
\begin{equation}
    \label{eq:population-maximization-problem}
    \max_{(N_A,N_B)\in\mathcal D} Y(N_A,N_B).
\end{equation}
This is not a dual problem. It is the same equilibrium-yield problem written
with different variables: after solving \eqref{eq:population-maximization-problem},
the corresponding efforts are recovered from
\eqref{eq:effort-recovery}.

The objective in \eqref{eq:population-yield} is concave. Its Hessian is
\begin{equation}
    \label{eq:yield-hessian}
    \nabla^2Y(N_A,N_B)
    =
    \begin{pmatrix}
        -\dfrac{2K_A^2r_A}{(K_A+N_A)^3} & 0\\
        0 & -\dfrac{2K_B^2r_B}{(K_B+N_B)^3}
    \end{pmatrix},
\end{equation}
Thus
\begin{equation}
    \label{eq:yield-hessian-negative-semidefinite}
    \nabla^2Y(N_A,N_B)\preceq0,
\end{equation}
so the Hessian is negative semidefinite for $K_A,K_B>0$ and
$r_A,r_B\geq0$.

We also need to check that the feasible population domain is convex. The
four boundary conditions in \eqref{eq:population-domain} can be
written in KKT form as
\begin{equation}
    \label{eq:population-constraints-kkt-form}
    \begin{aligned}
        -N_A&\leq0,
        &
        -N_B&\leq0,
        \\
        N_A-G_A(N_A,N_B)&\leq0,
        &
        N_B-G_B(N_A,N_B)&\leq0.
    \end{aligned}
\end{equation}
The first two inequalities are just $N_A\geq0$ and $N_B\geq0$ written in
the form $g_i\leq0$. Each one defines a half-plane, and half-planes are
convex.

For the last two inequalities, note that $G_A$ and $G_B$ are concave
functions of $(N_A,N_B)$. Indeed, since $0\leq m\leq1$, the coefficients
$m$ and $1-m$ are nonnegative. Therefore $G_A=(1-m)F_A+mF_B$ and
$G_B=mF_A+(1-m)F_B$ are nonnegative linear combinations of the concave
Beverton--Holt growth functions $F_A$ and $F_B$. Hence $G_A$ and $G_B$
are concave. Therefore $-G_A$ and $-G_B$ are convex. Adding the linear
terms $N_A$ and $N_B$ preserves convexity, so
\begin{equation*}
    N_A-G_A(N_A,N_B),
    \qquad
    N_B-G_B(N_A,N_B)
\end{equation*}
are convex functions. A sublevel set of a convex function,
$\{x:g_i(x)\leq0\}$, is convex
\cite[Section~3.1.6]{boyd-vandenberghe-2004}. Hence all four inequalities in
\eqref{eq:population-constraints-kkt-form} define convex sets.
The intersection of convex sets is convex
\cite[Section~2.3.1]{boyd-vandenberghe-2004}. Their intersection is exactly
$\mathcal D$, so $\mathcal D$ is convex.

Therefore the problem has a concave objective on a convex feasible domain.
For later use, label the four constraint functions as
\begin{equation}
    \label{eq:population-constraint-functions}
    \begin{aligned}
    g_1(N_A,N_B)&=-N_A,\\
    g_2(N_A,N_B)&=-N_B,\\
    g_3(N_A,N_B)&=N_A-G_A(N_A,N_B),\\
    g_4(N_A,N_B)&=N_B-G_B(N_A,N_B).
    \end{aligned}
\end{equation}
With multipliers $\lambda_1,\lambda_2,\lambda_3,\lambda_4\geq0$, the
Lagrangian for the population-domain yield problem is
\begin{equation}
    \label{eq:population-lagrangian}
    \begin{aligned}
    \mathcal L
    &=
    Y(N_A,N_B)
    -\lambda_1 g_1(N_A,N_B)
    -\lambda_2 g_2(N_A,N_B)\\
    &\qquad
    -\lambda_3 g_3(N_A,N_B)
    -\lambda_4 g_4(N_A,N_B).
    \end{aligned}
\end{equation}
Equivalently,
\begin{equation}
    \label{eq:population-lagrangian-expanded}
    \begin{aligned}
    \mathcal L
    &=
    Y(N_A,N_B)
    +\lambda_1N_A+\lambda_2N_B\\
    &\qquad
    -\lambda_3\bigl(N_A-G_A(N_A,N_B)\bigr)
    -\lambda_4\bigl(N_B-G_B(N_A,N_B)\bigr).
    \end{aligned}
\end{equation}
The cases below are active-set KKT certificates for
\eqref{eq:population-maximization-problem}: in each case we identify
the active constraints and verify
\eqref{eq:kkt-primal-feasibility}--\eqref{eq:kkt-stationarity}.

\subsubsection{Interior Case}

The interior case means that none of the four constraints in
\eqref{eq:population-constraints-kkt-form} is active. Equivalently,
\begin{equation*}
    N_A>0,\qquad N_B>0,\qquad N_A<G_A(N_A,N_B),\qquad
    N_B<G_B(N_A,N_B).
\end{equation*}
In effort variables this means
\begin{equation*}
    0<e_A<1,
    \qquad
    0<e_B<1.
\end{equation*}
Since all constraints are inactive, complementary slackness
\eqref{eq:kkt-complementary-slackness} is satisfied by taking all
multipliers equal to zero:
\begin{equation*}
    \lambda_1=\lambda_2=\lambda_3=\lambda_4=0.
\end{equation*}
Dual feasibility \eqref{eq:kkt-dual-feasibility} is then automatic,
and stationarity \eqref{eq:kkt-stationarity} reduces to the ordinary
first-order equations
\begin{equation}
    \label{eq:interior-stationarity}
    \frac{\partial Y}{\partial N_A}=0,
    \qquad
    \frac{\partial Y}{\partial N_B}=0.
\end{equation}
Using the separated yield formula \eqref{eq:population-yield}, the
partial derivatives are
\begin{equation}
    \label{eq:interior-yield-partials}
    \frac{\partial Y}{\partial N_A}=F_A'(N_A)-1,
    \qquad
    \frac{\partial Y}{\partial N_B}=F_B'(N_B)-1.
\end{equation}
Therefore the stationarity conditions
\eqref{eq:interior-stationarity} become
\begin{equation}
    \label{eq:interior-marginal-equations}
    F_A'(N_A^*)=1,
    \qquad
    F_B'(N_B^*)=1.
\end{equation}
For the Beverton--Holt growth functions,
\begin{equation}
    \label{eq:beverton-holt-derivatives}
    F_A'(N_A)=\frac{K_A^2r_A}{(K_A+N_A)^2},
    \qquad
    F_B'(N_B)=\frac{K_B^2r_B}{(K_B+N_B)^2}.
\end{equation}
Therefore the interior population candidate is
\begin{equation}
    \label{eq:interior-populations}
    N_A^*=K_A(\sqrt{r_A}-1),
    \qquad
    N_B^*=K_B(\sqrt{r_B}-1).
\end{equation}
This is an interior population candidate only when
\begin{equation}
    \label{eq:interior-r-condition}
    r_A>1,
    \qquad
    r_B>1.
\end{equation}

Substituting \eqref{eq:interior-populations} into the
Beverton--Holt maps gives
\begin{equation}
    \label{eq:interior-growth-values}
    F_A(N_A^*)=K_A\sqrt{r_A}(\sqrt{r_A}-1),
    \qquad
    F_B(N_B^*)=K_B\sqrt{r_B}(\sqrt{r_B}-1).
\end{equation}
Substituting \eqref{eq:interior-growth-values} into the definitions
of $G_A$ and $G_B$ gives
\begin{equation}
    \label{eq:interior-g-star}
    \begin{aligned}
    G_A^*
    &=
    (1-m)K_A\sqrt{r_A}(\sqrt{r_A}-1)
    +mK_B\sqrt{r_B}(\sqrt{r_B}-1),\\
    G_B^*
    &=
    mK_A\sqrt{r_A}(\sqrt{r_A}-1)
    +(1-m)K_B\sqrt{r_B}(\sqrt{r_B}-1).
    \end{aligned}
\end{equation}
The recovered efforts are
\begin{equation}
    \label{eq:interior-efforts}
    e_A^*
    =
    1-\frac{K_A(\sqrt{r_A}-1)}{G_A^*},
    \qquad
    e_B^*
    =
    1-\frac{K_B(\sqrt{r_B}-1)}{G_B^*}.
\end{equation}
When \eqref{eq:interior-r-condition} holds, we have
$N_A^*>0$, $N_B^*>0$, and $G_A^*,G_B^*>0$. Thus the remaining
conditions for the candidate to lie in the open interior of
$\mathcal D$ are
\begin{equation*}
    N_A^*<G_A^*,
    \qquad
    N_B^*<G_B^*.
\end{equation*}
Using the effort recovery formula \eqref{eq:effort-recovery}, these
two inequalities are equivalent to $e_A^*>0$ and $e_B^*>0$. Written in
parameters, they become
\begin{equation}
    \label{eq:interior-effort-validity}
    \scalebox{0.92}{$
    -K_B(\sqrt{r_B}-1)^2
    <
    m\bigl(K_A\sqrt{r_A}(\sqrt{r_A}-1)
    -K_B\sqrt{r_B}(\sqrt{r_B}-1)\bigr)
    <
    K_A(\sqrt{r_A}-1)^2.
    $}
\end{equation}
When equality holds in either side of
\eqref{eq:interior-effort-validity}, the candidate lies on a
no-take boundary rather than in the open interior.

The yield at the interior candidate is
\begin{equation}
    \label{eq:interior-yield}
    Y_{\mathrm{int}}
    =
    K_A(\sqrt{r_A}-1)^2
    +
    K_B(\sqrt{r_B}-1)^2.
\end{equation}

We can now check the KKT conditions. The inequalities
\eqref{eq:interior-r-condition} and
\eqref{eq:interior-effort-validity} give primal feasibility
\eqref{eq:kkt-primal-feasibility} with all four constraints strict.
All multipliers are zero, so dual feasibility
\eqref{eq:kkt-dual-feasibility} and complementary slackness
\eqref{eq:kkt-complementary-slackness} hold. Finally,
\eqref{eq:interior-marginal-equations} is exactly stationarity
\eqref{eq:kkt-stationarity}. Therefore, under
\eqref{eq:interior-r-condition} and
\eqref{eq:interior-effort-validity}, the interior candidate is a
global maximizer of the population-domain problem
\eqref{eq:population-maximization-problem}.

\begin{figure}[htbp]
    \centering
    \includegraphics[width=\textwidth]{plots/population_domain_view_interior.pdf}
    \caption{Interior example with $K_A=K_B=1$, $r_A=1.8$,
    $r_B=1.2$, and $m=0.2$. The left panel shows the yield over the
    original effort square. The right panel shows the same maximization
    problem in the population domain $\mathcal D$; the black curve is the
    boundary of the feasible population region. In both panels, the red
    point marks the analytic KKT candidate and the black cross marks the
    maximum found on the plotted grid.}
\end{figure}

\subsubsection{Complete-Harvest Boundaries}

We next consider the boundary where one patch has zero remaining
population after harvesting. Start with
\begin{equation*}
    N_A=0,
\end{equation*}
which corresponds to complete harvest in patch $A$, so $e_A=1$.

On this boundary, the feasibility condition for patch $B$ is
\begin{equation}
    \label{eq:complete-harvest-a-feasibility}
    0\leq N_B\leq G_B(0,N_B)=(1-m)F_B(N_B).
\end{equation}
For $N_B>0$, this inequality is equivalent to
\begin{equation*}
    K_B+N_B\leq K_B(1-m)r_B.
\end{equation*}
Thus the feasible interval on this boundary is
\begin{equation}
    \label{eq:complete-harvest-a-interval}
    0\leq N_B\leq K_B\bigl((1-m)r_B-1\bigr),
\end{equation}
provided $(1-m)r_B>1$. If $(1-m)r_B\leq1$, this boundary has no positive
feasible point.

The yield on this boundary is
\begin{equation}
    \label{eq:complete-harvest-a-yield-function}
    Y=F_B(N_B)-N_B.
\end{equation}
This is now a one-dimensional maximization problem on the interval
\eqref{eq:complete-harvest-a-interval}. A maximum on this interval
can occur either at an endpoint or inside the boundary interval. We first
test the second possibility, so assume
\begin{equation*}
    0<N_B<K_B\bigl((1-m)r_B-1\bigr).
\end{equation*}
At such a point, the endpoint constraints are inactive, and the
one-variable stationarity condition is
\begin{equation*}
    \frac{dY}{dN_B}=F_B'(N_B)-1=0.
\end{equation*}
Therefore the candidate on this boundary is
\begin{equation}
    \label{eq:complete-harvest-a-population}
    N_A^*=0,
    \qquad
    N_B^*=K_B(\sqrt{r_B}-1).
\end{equation}
Now we check whether this stationary point really lies inside the feasible
interval \eqref{eq:complete-harvest-a-interval}. This means
\begin{equation*}
    0<K_B(\sqrt{r_B}-1)<K_B\bigl((1-m)r_B-1\bigr).
\end{equation*}
The left inequality gives $r_B>1$. The right inequality gives
\begin{equation*}
    \sqrt{r_B}-1<(1-m)r_B-1
    \quad\Longleftrightarrow\quad
    \sqrt{r_B}<(1-m)r_B
    \quad\Longleftrightarrow\quad
    (1-m)\sqrt{r_B}>1,
\end{equation*}
where the last step uses $r_B>1$. Thus the candidate lies inside the
boundary interval exactly when
\begin{equation}
    \label{eq:complete-harvest-a-interval-condition}
    (1-m)\sqrt{r_B}>1.
\end{equation}
If equality holds, the boundary candidate lands exactly on the endpoint
$e_B=0$. The remaining possibility is
\begin{equation*}
    (1-m)r_B>1
    \qquad\text{but}\qquad
    (1-m)\sqrt{r_B}<1.
\end{equation*}
Then the feasible interval
\eqref{eq:complete-harvest-a-interval} contains positive values of
$N_B$, but the stationary point $K_B(\sqrt{r_B}-1)$ is to the right of the
interval. On the whole feasible interval we therefore have
\begin{equation*}
    N_B<K_B(\sqrt{r_B}-1).
\end{equation*}
Since $Y'(N_B)=F_B'(N_B)-1$ and $F_B'$ is decreasing, this means
\begin{equation*}
    Y'(N_B)>0
\end{equation*}
throughout the feasible interval. Thus $Y$ is increasing on the whole
boundary interval, so the constrained boundary maximum is attained at the
right endpoint. This right endpoint is the case $N_B=G_B(0,N_B)$, hence
$e_B=0$, and it is treated in the next section.

Now derive the KKT multiplier. At this candidate, the active constraint is
$g_1=-N_A=0$. The other constraints are inactive, so complementary
slackness gives
\begin{equation*}
    \lambda_2=\lambda_3=\lambda_4=0.
\end{equation*}
Stationarity in the $N_B$ direction is exactly
$F_B'(N_B^*)-1=0$. In the $N_A$ direction, the expanded Lagrangian
\eqref{eq:population-lagrangian-expanded} gives
\begin{equation*}
    0=\frac{\partial\mathcal L}{\partial N_A}
    =
    \bigl(F_A'(0)-1\bigr)+\lambda_1
    =
    (r_A-1)+\lambda_1.
\end{equation*}
Hence
\begin{equation}
    \label{eq:complete-harvest-a-multipliers}
    \lambda_1=1-r_A,
    \qquad
    \lambda_2=\lambda_3=\lambda_4=0.
\end{equation}
Now dual feasibility requires $\lambda_i\geq0$ for every multiplier. In
particular, $\lambda_1\geq0$, so
\begin{equation*}
    1-r_A\geq0.
\end{equation*}
Therefore
\begin{equation}
    \label{eq:complete-harvest-a-multiplier-condition}
    r_A\leq1.
\end{equation}

Combining \eqref{eq:complete-harvest-a-interval-condition} and
\eqref{eq:complete-harvest-a-multiplier-condition}, complete harvest
in patch $A$ gives a KKT certificate when
\begin{equation}
    \label{eq:complete-harvest-a-condition}
    r_A\leq1,
    \qquad
    (1-m)\sqrt{r_B}>1.
\end{equation}
The recovered efforts are
\begin{equation}
    \label{eq:complete-harvest-a-efforts}
    e_A^*=1,
    \qquad
    e_B^*=1-\frac{1}{(1-m)\sqrt{r_B}},
\end{equation}
and the yield is
\begin{equation}
    \label{eq:complete-harvest-a-yield}
    Y^*=K_B(\sqrt{r_B}-1)^2.
\end{equation}
Primal feasibility follows from the boundary interval, dual feasibility is
\eqref{eq:complete-harvest-a-multiplier-condition}, complementary
slackness holds because inactive constraints have zero multipliers, and
stationarity was checked above. Therefore, when
\eqref{eq:complete-harvest-a-condition} holds, this boundary
candidate is globally optimal.

The complete-harvest-in-$B$ case is symmetric. The active boundary is
$N_B=0$, equivalently $e_B=1$. If
\begin{equation}
    \label{eq:complete-harvest-b-condition}
    r_B\leq1,
    \qquad
    (1-m)\sqrt{r_A}>1,
\end{equation}
then the KKT certificate gives
\begin{equation}
    \label{eq:complete-harvest-b-population}
    N_B^*=0,
    \qquad
    N_A^*=K_A(\sqrt{r_A}-1),
\end{equation}
with efforts
\begin{equation}
    \label{eq:complete-harvest-b-efforts}
    e_B^*=1,
    \qquad
    e_A^*=1-\frac{1}{(1-m)\sqrt{r_A}},
\end{equation}
and yield
\begin{equation}
    \label{eq:complete-harvest-b-yield}
    Y^*=K_A(\sqrt{r_A}-1)^2.
\end{equation}

\begin{figure}[htbp]
    \centering
    \includegraphics[width=\textwidth]{plots/population_domain_view_complete_harvest.pdf}
    \caption{Complete-harvest examples. Each row shows the same parameter
    set in the effort domain on the left and in the population domain on
    the right. Complete harvest appears as $e_A=1$ or $e_B=1$ in the
    effort domain, and as $N_A=0$ or $N_B=0$ in the population domain.}
\end{figure}

\subsubsection{Protect Source, Harvest Sink Endpoint}

The previous boundary calculation also points to an important endpoint.
On the boundary $N_A=0$, the right endpoint of the feasible interval
\eqref{eq:complete-harvest-a-interval} is
\begin{equation}
    \label{eq:source-endpoint-a-sink-population}
    N_A^*=0,
    \qquad
    N_B^*=K_B\bigl((1-m)r_B-1\bigr).
\end{equation}
This point exists with $N_B^*>0$ when
\begin{equation}
    \label{eq:source-endpoint-a-sink-positive}
    (1-m)r_B>1.
\end{equation}
At this endpoint,
\begin{equation*}
    N_B^*=G_B(0,N_B^*),
\end{equation*}
so the recovered efforts are
\begin{equation}
    \label{eq:source-endpoint-a-sink-efforts}
    (e_A^*,e_B^*)=(1,0).
\end{equation}
Patch $B$ is protected and acts as the source, while patch $A$ is
harvested completely as the sink.

The yield is especially simple. Since
$N_B^*=(1-m)F_B(N_B^*)$, we have
\begin{equation*}
    F_B(N_B^*)=\frac{N_B^*}{1-m}.
\end{equation*}
Therefore
\begin{equation}
    \label{eq:source-endpoint-a-sink-yield}
    Y^*
    =
    F_B(N_B^*)-N_B^*
    =
    \frac{m}{1-m}
    K_B\bigl((1-m)r_B-1\bigr).
\end{equation}

Now check the KKT multipliers. The active constraints are
\begin{equation*}
    g_1=-N_A=0,
    \qquad
    g_4=N_B-G_B(N_A,N_B)=0.
\end{equation*}
The other two constraints are inactive, so
$\lambda_2=\lambda_3=0$. Using
\eqref{eq:source-endpoint-a-sink-population}, the marginal growth in
patch $B$ at this endpoint is
\begin{equation}
    \label{eq:source-endpoint-a-sink-derivative-b}
    F_B'(N_B^*)
    =
    \frac{1}{(1-m)^2r_B}.
\end{equation}
Stationarity in the $N_B$ direction gives
\begin{equation*}
    0
    =
    \left(\frac{1}{(1-m)^2r_B}-1\right)
    -
    \lambda_4
    \left(1-\frac{1}{(1-m)r_B}\right),
\end{equation*}
and therefore
\begin{equation}
    \label{eq:source-endpoint-a-sink-lambda-b}
    \lambda_4
    =
    \frac{1-(1-m)^2r_B}
         {(1-m)\bigl((1-m)r_B-1\bigr)}.
\end{equation}
Now check dual feasibility for this multiplier. Since
\eqref{eq:source-endpoint-a-sink-positive} gives
$(1-m)r_B-1>0$, the denominator in
\eqref{eq:source-endpoint-a-sink-lambda-b} is positive. Therefore
$\lambda_4\geq0$ is equivalent to the numerator being nonnegative:
\begin{equation*}
    1-(1-m)^2r_B\geq0.
\end{equation*}
Equivalently,
\begin{equation}
    \label{eq:source-endpoint-a-sink-lambda-b-condition}
    (1-m)\sqrt{r_B}\leq1.
\end{equation}

It remains to check the multiplier for $g_1=-N_A=0$. Stationarity in the
$N_A$ direction gives
\begin{equation*}
    0
    =
    (r_A-1)+\lambda_1+m r_A\lambda_4.
\end{equation*}
Thus
\begin{equation}
    \label{eq:source-endpoint-a-sink-lambda-a}
    \lambda_1
    =
    1-r_A(1+m\lambda_4).
\end{equation}
Dual feasibility for $\lambda_1$ requires
\begin{equation}
    \label{eq:source-endpoint-a-sink-ra-condition}
    r_A\leq \frac{1}{1+m\lambda_4},
\end{equation}
with $\lambda_4$ given by
\eqref{eq:source-endpoint-a-sink-lambda-b}.
Substituting \eqref{eq:source-endpoint-a-sink-lambda-b}, this can
be written explicitly as
\begin{equation}
    \label{eq:source-endpoint-a-sink-ra-condition-explicit}
    r_A
    \leq
    \frac{(1-m)\bigl((1-m)r_B-1\bigr)}
         {(1-m)^3r_B+2m-1}.
\end{equation}

Combining the conditions, the protect-$B$, harvest-$A$ endpoint gives a
KKT certificate when
\begin{equation}
    \label{eq:source-endpoint-a-sink-condition}
    (1-m)r_B>1,
    \qquad
    (1-m)\sqrt{r_B}\leq1,
    \qquad
    r_A
    \leq
    \frac{(1-m)\bigl((1-m)r_B-1\bigr)}
         {(1-m)^3r_B+2m-1}.
\end{equation}
Under these conditions, primal feasibility follows from the endpoint
construction, dual feasibility follows from
\eqref{eq:source-endpoint-a-sink-lambda-b-condition} and
\eqref{eq:source-endpoint-a-sink-ra-condition}, complementary
slackness holds by the choice of active constraints, and stationarity was
derived above. Therefore this endpoint is globally optimal.

The symmetric endpoint protects $A$ and harvests $B$, so
\begin{equation}
    \label{eq:source-endpoint-b-sink-population}
    N_A^*=K_A\bigl((1-m)r_A-1\bigr),
    \qquad
    N_B^*=0,
\end{equation}
with efforts
\begin{equation}
    \label{eq:source-endpoint-b-sink-efforts}
    (e_A^*,e_B^*)=(0,1).
\end{equation}
Here
\begin{equation}
    \label{eq:source-endpoint-b-sink-lambda-a}
    \lambda_3
    =
    \frac{1-(1-m)^2r_A}
         {(1-m)\bigl((1-m)r_A-1\bigr)}.
\end{equation}
This symmetric endpoint gives a KKT certificate when
\begin{equation}
    \label{eq:source-endpoint-b-sink-condition}
    (1-m)r_A>1,
    \qquad
    (1-m)\sqrt{r_A}\leq1,
    \qquad
    r_B
    \leq
    \frac{(1-m)\bigl((1-m)r_A-1\bigr)}
         {(1-m)^3r_A+2m-1}.
\end{equation}
The yield is
\begin{equation}
    \label{eq:source-endpoint-b-sink-yield}
    Y^*
    =
    \frac{m}{1-m}
    K_A\bigl((1-m)r_A-1\bigr).
\end{equation}

\begin{figure}[htbp]
    \centering
    \includegraphics[width=\textwidth]{plots/population_domain_view_source_endpoint.pdf}
    \caption{Protect-source, harvest-sink endpoint examples. In effort
    variables these are the corner strategies $(e_A,e_B)=(0,1)$ and
    $(1,0)$. In population variables the same endpoints lie where one
    patch has zero remaining population and the other lies on its no-take
    boundary.}
\end{figure}

\subsubsection{No-Take Boundaries}

One patch is protected from fishing, while the other patch may still be
harvested partially. This is an important case, since a common
conservation strategy is to declare a marine reserve where fishing is not
allowed. Thus it is useful to study the optimal fishing rate in the other
patch, even when the no-take strategy is not globally optimal. It is also
important to know when declaring a marine reserve is itself the optimal
choice.

We first consider no take in patch $A$. In population variables this means
\begin{equation}
    \label{eq:no-take-a-active-constraint}
    N_A=G_A(N_A,N_B),
\end{equation}
so the active constraint is $g_3=0$. For a nonendpoint no-take solution,
the other constraints are inactive:
\begin{equation*}
    N_A>0,\qquad N_B>0,\qquad N_B<G_B(N_A,N_B).
\end{equation*}
Thus complementary slackness gives
\begin{equation*}
    \lambda_1=\lambda_2=\lambda_4=0,
\end{equation*}
and only $\lambda_3$ remains.

From the expanded Lagrangian
\eqref{eq:population-lagrangian-expanded}, stationarity gives
\begin{equation}
    \label{eq:no-take-a-kkt-stationarity}
    \begin{aligned}
    1-F_A'(N_A)
    +\lambda_3\bigl(1-(1-m)F_A'(N_A)\bigr)&=0,\\
    1-F_B'(N_B)-m\lambda_3F_B'(N_B)&=0.
    \end{aligned}
\end{equation}
The second equation gives
\begin{equation}
    \label{eq:no-take-a-multiplier}
    \lambda_3=\frac{1-F_B'(N_B)}{mF_B'(N_B)}
\end{equation}
when $m>0$. Therefore dual feasibility $\lambda_3\geq0$ is equivalent to
\begin{equation}
    \label{eq:no-take-a-dual-feasibility}
    F_B'(N_B)\leq1.
\end{equation}
Substituting \eqref{eq:no-take-a-multiplier} into the first
stationarity equation and clearing the denominator gives the no-take
stationarity relation
\begin{equation}
    \label{eq:no-take-stationarity-derivatives}
    (1-m)\bigl(F_A'(N_A)+F_B'(N_B)\bigr)
    +(2m-1)F_A'(N_A)F_B'(N_B)=1.
\end{equation}

The stationarity relation can be solved directly for the marginal growth
in patch $B$:
\begin{equation}
    \label{eq:no-take-a-derivative-b-from-a}
    F_B'(N_B)
    =
    \frac{1-(1-m)F_A'(N_A)}
         {1-m+(2m-1)F_A'(N_A)}.
\end{equation}
For the Beverton--Holt map,
\begin{equation*}
    F_B'(N_B)=\frac{K_B^2r_B}{(K_B+N_B)^2}.
\end{equation*}
Combining this with \eqref{eq:no-take-a-derivative-b-from-a} gives
\begin{equation*}
    \frac{K_B^2r_B}{(K_B+N_B)^2}
    =
    \frac{1-(1-m)F_A'(N_A)}
         {1-m+(2m-1)F_A'(N_A)}.
\end{equation*}
Taking the positive square root gives
\begin{equation*}
    K_B+N_B
    =
    K_B
    \sqrt{
        \frac{
        r_B\bigl(1-m+(2m-1)F_A'(N_A)\bigr)}
        {1-(1-m)F_A'(N_A)}
    }.
\end{equation*}
Thus
\begin{equation}
    \label{eq:no-take-a-nb-from-na}
    \widehat N_B(N_A)
    =
    K_B\left(
    \sqrt{
        \frac{
        r_B\bigl(1-m+(2m-1)F_A'(N_A)\bigr)}
        {1-(1-m)F_A'(N_A)}
    }
    -1\right).
\end{equation}
The stationarity condition alone does not determine the equilibrium. The
protected patch equation \eqref{eq:no-take-a-active-constraint} must
also hold. Substituting $N_B=\widehat N_B(N_A)$ gives the scalar equation
\begin{equation}
    \label{eq:no-take-a-scalar-root}
    N_A
    =
    (1-m)F_A(N_A)+mF_B\bigl(\widehat N_B(N_A)\bigr).
\end{equation}
After solving \eqref{eq:no-take-a-scalar-root} for $N_A$, set
$N_B=\widehat N_B(N_A)$ and recover the effort in patch $B$ from
\eqref{eq:effort-recovery}:
\begin{equation}
    \label{eq:no-take-a-effort-recovery}
    e_B^*
    =
    1-
    \frac{N_B}
         {mF_A(N_A)+(1-m)F_B(N_B)}.
\end{equation}

For this active set the remaining primal checks are
\begin{equation}
    \label{eq:no-take-a-primal-validity}
    N_A>0,\qquad N_B>0,\qquad N_B<G_B(N_A,N_B).
\end{equation}
By \eqref{eq:effort-recovery}, the last inequality is equivalent to
$e_B^*>0$, so the recovered effort lies on the open harvested edge. The
dual feasibility check is
\begin{equation}
    \label{eq:no-take-a-dual-validity}
    F_B'(N_B)\leq1.
\end{equation}
Under these conditions, primal feasibility follows from
\eqref{eq:no-take-a-active-constraint} and
\eqref{eq:no-take-a-primal-validity}. For dual feasibility, the
only nonzero multiplier is $\lambda_3$, and
\eqref{eq:no-take-a-multiplier} shows that
\eqref{eq:no-take-a-dual-validity} is exactly the condition
$\lambda_3\geq0$.
Complementary slackness holds because only $g_3$ is active, and
stationarity is exactly
\eqref{eq:no-take-a-kkt-stationarity}. Therefore this no-take
candidate is globally optimal.

The no-take-in-$B$ boundary is symmetric. Here $g_4=0$, or
\begin{equation}
    \label{eq:no-take-b-active-constraint}
    N_B=G_B(N_A,N_B).
\end{equation}
The same stationarity relation
\eqref{eq:no-take-stationarity-derivatives} gives
\begin{equation}
    \label{eq:no-take-b-derivative-a-from-b}
    F_A'(N_A)
    =
    \frac{1-(1-m)F_B'(N_B)}
         {1-m+(2m-1)F_B'(N_B)}.
\end{equation}
Thus
\begin{equation}
    \label{eq:no-take-b-na-from-nb}
    \widehat N_A(N_B)
    =
    K_A\left(
    \sqrt{
        \frac{
        r_A\bigl(1-m+(2m-1)F_B'(N_B)\bigr)}
        {1-(1-m)F_B'(N_B)}
    }
    -1\right).
\end{equation}
The protected patch equation becomes the scalar equation
\begin{equation}
    \label{eq:no-take-b-scalar-root}
    N_B
    =
    mF_A\bigl(\widehat N_A(N_B)\bigr)+(1-m)F_B(N_B).
\end{equation}
After solving \eqref{eq:no-take-b-scalar-root} for $N_B$, set
$N_A=\widehat N_A(N_B)$ and recover the effort in patch $A$:
\begin{equation}
    \label{eq:no-take-b-effort-recovery}
    e_A^*
    =
    1-
    \frac{N_A}
         {(1-m)F_A(N_A)+mF_B(N_B)}.
\end{equation}
The symmetric primal checks are
\begin{equation}
    \label{eq:no-take-b-primal-validity}
    N_A>0,\qquad N_B>0,\qquad N_A<G_A(N_A,N_B).
\end{equation}
By \eqref{eq:effort-recovery}, the last inequality is equivalent to
$e_A^*>0$. The symmetric dual feasibility check is
\begin{equation}
    \label{eq:no-take-b-dual-validity}
    F_A'(N_A)\leq1.
\end{equation}
Here the only nonzero multiplier is $\lambda_4$, and the symmetric
multiplier formula is
\begin{equation*}
    \lambda_4=\frac{1-F_A'(N_A)}{mF_A'(N_A)}.
\end{equation*}
Thus \eqref{eq:no-take-b-dual-validity} is exactly the condition
$\lambda_4\geq0$. When \eqref{eq:no-take-b-primal-validity} and
\eqref{eq:no-take-b-dual-validity} hold, the no-take-in-$B$
candidate is globally optimal.

\begin{figure}[htbp]
    \centering
    \includegraphics[width=\textwidth]{plots/population_domain_view_no_take.pdf}
    \caption{No-take examples. The effort-domain view shows the protected
    patch as $e_A=0$ or $e_B=0$. The population-domain view shows the
    corresponding active constraint $N_A=G_A(N_A,N_B)$ or
    $N_B=G_B(N_A,N_B)$ as part of the feasible-region boundary.}
\end{figure}

\subsubsection{KKT Maximizer Certificates}

First check whether the linked population can persist without fishing. If
it cannot, then there is no positive sustainable yield to maximize. After
that extinction check, the rows below are independent KKT certificates.
Whenever the conditions in a row hold, the corresponding candidate is a
global maximizer of \eqref{eq:population-maximization-problem}; the
rows do not have to be checked in this order. The no-take rows are
implicit conditions: to check them, one must first solve the corresponding
scalar equation and then verify the primal and dual checks for that
solution.

\begin{center}
\small
\setlength{\tabcolsep}{3pt}
\renewcommand{\arraystretch}{1.35}
\begin{tabular}{
    >{\raggedright\arraybackslash}p{0.13\textwidth}
    >{\raggedright\arraybackslash}p{0.29\textwidth}
    >{\raggedright\arraybackslash}p{0.30\textwidth}
    >{\raggedright\arraybackslash}p{0.17\textwidth}
}
\hline
Case & If this condition holds & Maximizing efforts & Yield \\
\hline
Interior
&
\eqref{eq:interior-r-condition} and
\eqref{eq:interior-effort-validity}
&
\eqref{eq:interior-efforts}
&
\eqref{eq:interior-yield}
\\

Complete harvest in $A$
&
\eqref{eq:complete-harvest-a-condition}
&
\eqref{eq:complete-harvest-a-efforts}
&
\eqref{eq:complete-harvest-a-yield}
\\

Complete harvest in $B$
&
\eqref{eq:complete-harvest-b-condition}
&
\eqref{eq:complete-harvest-b-efforts}
&
\eqref{eq:complete-harvest-b-yield}
\\

Protect $B$, harvest $A$
&
\eqref{eq:source-endpoint-a-sink-condition}
&
\eqref{eq:source-endpoint-a-sink-efforts}
&
\eqref{eq:source-endpoint-a-sink-yield}
\\

Protect $A$, harvest $B$
&
\eqref{eq:source-endpoint-b-sink-condition}
&
\eqref{eq:source-endpoint-b-sink-efforts}
&
\eqref{eq:source-endpoint-b-sink-yield}
\\

No take in $A$
&
The solution of \eqref{eq:no-take-a-scalar-root} satisfies
\eqref{eq:no-take-a-primal-validity},
\eqref{eq:no-take-a-dual-validity}.
&
$e_A^*=0$, and $e_B^*$ from
\eqref{eq:no-take-a-effort-recovery}
&
$G_B(N_A,N_B)-N_B$
\\

No take in $B$
&
The solution of \eqref{eq:no-take-b-scalar-root} satisfies
\eqref{eq:no-take-b-primal-validity},
\eqref{eq:no-take-b-dual-validity}.
&
$e_B^*=0$, and $e_A^*$ from
\eqref{eq:no-take-b-effort-recovery}
&
$G_A(N_A,N_B)-N_A$
\\
\hline
\end{tabular}
\end{center}

Figure~\ref{fig:numeric-zone-map} gives a numerical check of these
certificates over a rectangle of parameter values. For each pair
$(r_A,r_B)$, the color shows which candidate family gives the largest
simulated yield. The simulation first searches a coarse grid in
$(e_A,e_B)$, then refines only the coarse grid cell containing the best
point together with its neighboring cells. The dashed and dotted curves
are the analytic boundaries from the conditions in the table.

\begin{figure}[H]
    \centering
    \includegraphics[width=0.9\textwidth]{plots/numeric_yield_zone_map_ka2_kb1_m0_45.pdf}
    \caption{Numerical region map for the yield-maximizing candidate
    family with $m=0.45$, $K_A=2$, and $K_B=1$. The colored background is
    produced by direct simulation in effort variables. The overlaid
    curves are the analytic boundaries from the KKT certificates above.}
    \label{fig:numeric-zone-map}
\end{figure}

\section{Conclusion}

This paper started with the one-population Beverton--Holt model and then
added harvesting through a fishing-effort parameter. In that setting the
equilibria, extinction threshold, and yield-maximizing effort can be
computed explicitly. The same questions become richer in the linked
two-patch model, because dispersal makes the fate of each patch depend on
the other patch.
\\
For the linked model, we first studied extinction. Linearization at the
extinction equilibrium gives a matrix condition, and the spectral radius
of this matrix determines when extinction is locally stable. This gives a
clear way to compare the analytic extinction boundary with simulated
dynamics.
\\
We then studied yield maximization for the linked model with fishing.
The main simplification was to reformulate the problem in the equilibrium
population variables rather than directly in the fishing efforts. In
these variables the feasible set is convex and the yield function is
concave, so the KKT conditions provide global optimality certificates.
This led to several possible optimal strategies: an interior fishing
strategy, complete harvest in one patch, protection of a source patch
while harvesting a sink patch, and no-take marine reserve boundaries.
The numerical region map shows how these cases divide parameter space.
\\
One natural future direction is to optimize the connectivity itself. In
this paper the movement parameter $m$ is fixed. In some real-world
settings, management may also influence connectivity: for example, by
opening or closing migration corridors, changing habitat quality between
regions, or modifying barriers that affect movement. Then the yield
maximization problem would have three decision variables,
\begin{equation*}
    e_A,\qquad e_B,\qquad m.
\end{equation*}
The same general method may still be useful: write the feasible
population domain, recover the efforts from the equilibrium constraints,
and apply KKT active-set analysis. This extension is outside the scope of
the present paper, where the goal was to understand the effect of fishing
efforts for a fixed linked system.

\begin{thebibliography}{9}

\bibitem{beverton-holt-1957}
R. J. H. Beverton and S. J. Holt.
\textit{On the Dynamics of Exploited Fish Populations}.
Fishery Investigations, Series II, Volume XIX.
Ministry of Agriculture, Fisheries and Food, London, 1957.

\bibitem{kinderlehrer-stampacchia-1980}
D. Kinderlehrer and G. Stampacchia.
Variational inequalities in $R^N$.
In \textit{An Introduction to Variational Inequalities and Their
Applications}, pages 7--22.
Academic Press, New York, 1980.

\bibitem{boyd-vandenberghe-2004}
S. Boyd and L. Vandenberghe.
\textit{Convex Optimization}.
Cambridge University Press, Cambridge, 2004.

\bibitem{elaydi-2005-stability}
S. N. Elaydi.
Stability theory.
In \textit{An Introduction to Difference Equations}, third edition,
pages 173--243.
Springer, New York, 2005.
doi: 10.1007/0-387-27602-5\_4.

\end{thebibliography}
\end{document}
